\documentclass[11pt,a4paper]{letter}

\usepackage[margin=0.9in,top=0.7in,bottom=0.7in]{geometry}
\usepackage{hyperref}
\renewcommand{\baselinestretch}{0.97}

\signature{Liang Wang}
\address{School of Artificial Intelligence and Automation \\
Huazhong University of Science and Technology \\
Wuhan 430074, P.R. China \\
\texttt{wangliang.f@gmail.com}}

\begin{document}

\begin{letter}{
The Editors \\
\emph{Nonlinearity} \\
IOP Publishing \& London Mathematical Society
}

\opening{Dear Editors,}

I am submitting the enclosed manuscript, \emph{Transient Chaos and
Topological Bounds in Prime Dynamics: Revisiting the One-Dimensional
Sieve Mapping}, for consideration as a research article.

The paper studies the symbolic dynamics of the Eratosthenes sieve
under the Metropolis--Stein--Stein (MSS) kneading admissibility
framework, taking as starting object the matching, recently introduced
in a published companion paper (Wang, \emph{Research in Mathematics}
\textbf{13}(1), 2026, DOI: \href{https://doi.org/10.1080/27684830.2026.2684334}{10.1080/27684830.2026.2684334}),
between the sieve word $Q_k$ and admissible itineraries of the
unimodal map $x \mapsto 1 - u x^2$ at the band-merging parameter
$u_c \approx 1.5436890$. Two unconditional theorems anchor the
contribution. The \emph{Parity-Gap Lemma} (Theorem~A) shows that any
MSS-defective shift of $W_k = Q_k[0, p_{k+1}^2)$ forces a prime gap
of length $\geq p_{k+1} - 1$ at a specific position; this reduces
the topological-admissibility question to an extremal prime-gap
inequality, and we observe that classical gap bounds (Legendre,
Andrica, RH, BHP) all fail to clear it---only a strongly sub-root
bound such as Cram\'er's conjecture forces eventual admissibility.
The \emph{even-gap rigidity} theorem (Theorem~B) shows that gaps
between consecutive $L$-symbols at $u_c$ are almost surely even,
via the parity-twisted lex order on $RLR^\infty$. A third structural
result (Theorem~C) shows the gap measure decays asymptotically
geometrically with ratio $p_\infty \approx 0.596$, ruling out
internal mod-$3$ resonance and identifying the unimodal projection
as an abelian / mod-$2$ shadow of the prime distribution. The paper
closes with a clearly conditional reduction connecting recurrence
to the $L$-$R$-$L$ cylinder to a precisely stated arithmetic
shadowing problem; we explicitly disclaim any consequence for
classical number-theoretic conjectures.

The paper fits squarely within \emph{Nonlinearity}'s remit at the
intersection of symbolic dynamics, ergodic theory of one-dimensional
unimodal maps, and arithmetic dynamics. The proofs rely on classical
results of Misiurewicz, Ma\~n\'e, Lasota--Yorke, and Young, which are
familiar to the journal's audience. The methodological pattern of
the paper---numerical experiments to surface a structural lemma
followed by a proof on classical foundations---is in keeping with
\emph{Nonlinearity}'s tradition of bridging quantitative experiment
and rigorous theory.

I should disclose that I am also the author of the published
companion paper cited above, which proposes the unimodal--sieve
matching itself; the present paper independently analyses the
admissibility structure of that matching and is structured to stand
alone---no theorem statement depends on the prior work, as recorded
in the title-page footnote. The manuscript is not under consideration
elsewhere; there are no competing interests; and the complete LaTeX
source, the numerical experiments, the figure-generation scripts, and
six companion Jupyter notebooks reproducing the headline numerical
claims are openly available at
\url{https://github.com/maris205/prime_dynamics}.

I thank the Editors and referees in advance for their consideration.

\closing{Sincerely yours,}

\end{letter}

\end{document}
